Investiga-se a confiabilidade da inferência da massa do gráviton com redes
de temporização de pulsares (PTAs), comparando frequência explícita,
compressão consistente e resposta avaliada em uma frequência de referência.
Um controle adicional conserva os coeficientes Fourier. A metodologia
separa perda de informação, erro de resposta e aproximação probabilística.
O estudo parte do TG de Wayne Leonardo Silva de Paula (ITA, 2003), publicado
em 2004; a bibliografia foi atualizada até 11 de setembro de 2026.

\noindent\textbf{Alcance da versão \VersaoTrabalho.} Foram concluídos os
fundamentos, a resposta tensorial, o simulador e as comparações de inferência
e compressão. Na campanha de 500 realizações, as aproximações normais nas
estatísticas físicas tiveram dez rejeições persistentes; 243 PITs de 15000
permaneceram não resolvidos. A comparação pareada de nove análises sobre
32 dados é inconclusiva na escala de 0,01 da largura da priori após propagar
a incerteza numérica. Não se demonstra equivalência de limites de massa.

\noindent\textbf{Robustez condicional.} C09 reúne 25\,956 posteriores de
base, um painel de 1\,600 integrações de prioris e 39 contrastes pareados.
A diagonalização tem efeitos de ambos os sinais; omitir contaminantes pode
reduzir artificialmente limites. A família SBC de 18 grupos e 126 testes
apresenta 116 não rejeições condicionais e dez resultados numericamente
indeterminados. Permanecem 834 PITs de massa e 2\,287 eventos sem resolução
operacional na produção completa. Os domínios interpolados não possuem
certificado físico uniforme; não rejeitar não demonstra aprendizado da massa.
O capítulo conserva o histórico de falhas e distingue a produção dos pilotos.

\noindent\textbf{Capítulos em draft.} A campanha escalar bidimensional e a
aplicação observacional de C10 e C11 permanecem pendentes. C10--C13 continuam
abertos; não se apresenta nova restrição observacional ou detecção.
